% =====================================================================
% REWRITTEN SECTIONS for the STRIPE paper, grounded in the actual pipeline
% (config_coherent_power_perf_perfreq_single_channel.yaml and the operator graph).
%
% PART A  -> additions to \section{End-to-End System Architecture}:
%             (A1) a real-time framing paragraph to open the section, and
%             (A2) a new \subsection on offline replay + loopback validation.
%           The existing ingest/FFT/detector/visualization subsections are
%           accurate as written; keep them (see the accuracy notes in chat).
%
% PART B  -> a single \section{Detector Variants} that REPLACES both the old
%           \section{Coherent-Power Detector} (Sec.~IV) and the old
%           \section{Detector Variants} (Sec.~V); coherent power is folded in
%           as its real-time-default subsection.
% =====================================================================


% =====================================================================
% PART A1 -- opening paragraph for \section{End-to-End System Architecture}
% (insert immediately after \label{sec:system}, before the Fig. reference)
% =====================================================================

STRIPE is designed to run \emph{at line rate} on the streaming data path: the
detector consumes complex FFT frames as they are produced, and no stage stages the
raw sample firehose through host memory or disk. Our reference real-time
configuration is the single-channel, calibrated per-frequency coherent-power profile
(\texttt{config\_coherent\_power\_perf\_perfreq\_single\_channel.yaml}), which
sustains a \SI{500}{MS/s} complex stream end-to-end on a single GPU. Multi-channel
operation instantiates one replicated graph per channel, unless a detector variant
explicitly consumes cross-channel features. The active operator chain is CHDR
ingest $\rightarrow$ wideband FFT $\rightarrow$ detector, with mask preview,
spectrogram preview, and logging attached as non-blocking downstream consumers;
the display spectrogram is a consumer-side product, not an intermediate on the hot
path.


% =====================================================================
% PART A2 -- NEW subsection (append at the end of \section{End-to-End System Architecture})
% =====================================================================

\subsection{Offline Replay and Loopback Validation}
\label{sec:replay}

Because the same operator graph is used online and offline, STRIPE supports two
complementary replay modes for controlled, reproducible experiments.

\emph{(i) File replay for detector evaluation.} A stored SigMF capture is streamed
through the graph in software: the FFT geometry is derived from the file's sample
rate, frames are reconstructed identically to the live path, and any detector
operator is run over them at sub-real-time pace. Because the FFT, detector, and
scoring code are shared with the deployed graph, this isolates the \emph{detector}
as the only variable, so competing detectors are compared on byte-identical input
against SigMF ground truth. All quantitative results in Sec.~\ref{sec:eval} use this
mode.

\emph{(ii) Cable loopback for full-system validation.} A known waveform is
transmitted from the USRP and received back through the complete RF front end and
the \emph{same} DPDK/GPUDirect CHDR ingest path used in deployment, then processed
by the live FFT and detector operators. Unlike file replay, this exercises the
entire end-to-end pathway---packetization, GPU DMA, batch assembly, timing and
frequency metadata, and soft-resync recovery---on a reproducible, physically
transmitted signal (the ingest stage merely relaxes its partial-batch timeout to
tolerate replay pacing). Loopback thus validates that the online system reproduces
the behavior measured offline, closing the loop between controlled evaluation and
real-time deployment.


% =====================================================================
% PART B -- REPLACEMENT for old Sec.~IV (Coherent-Power Detector) AND
%           old Sec.~V (Detector Variants).
% =====================================================================

\section{Detector Variants}
\label{sec:variants}

STRIPE is detector-agnostic: any operator that consumes a complex FFT frame (or a
derived feature tensor) and emits a full-resolution binary mask plugs into the graph
behind a common interface, so the detector becomes a runtime systems knob. We
implement and evaluate three families spanning the accuracy--compute tradeoff: a
deployed coherent-power detector as the low-cost real-time default, classical
signal-processing baselines, and learned (foundation-model and object-detector)
variants. Two families run as live GPU operators (\emph{coherent\_power} and
zero-shot DINO); the remaining variants are evaluated through the offline file-replay
path (Sec.~\ref{sec:replay}) on identical input, so all comparisons in
Sec.~\ref{sec:eval} are apples-to-apples.

\subsection{Coherent-Power Detector (Real-Time Default)}
\label{sec:coherent}

The deployed real-time detector is a hybrid image/power algorithm implemented as CUDA
kernels; it treats the dB-power spectrogram both as a 2-D image and as a bank of
calibrated per-frequency power channels, and emits the union of the two views after
morphological cleanup. From a complex frame $X[f,t]$ it forms
$P_{\db}=10\log_{10}(|X|^2+10^{-12})$; an optional per-frequency front-end
equalization (Gaussian row-mean smoothing, $\sigma=12$ bins) is available but clamped
to \SI{0}{dB} in the real-time profile.

\emph{Image view (local support).} A separable box filter estimates a local
background $B[f,t]$ over a $\pm10$-bin time and $\pm8$-bin frequency window; the
normalized excess $S=P_{\db}-B$ is thresholded (floor \SI{0.5}{dB}, span \SI{3}{dB},
$\tau=0.4$), admitting pixels roughly \SI{1.7}{dB} above their local neighborhood.
This CFAR-like test tracks slow gain variation but hollows the interior of signals
wider than the box.

\emph{Power view (fill + rescue).} Two absolute-power tests fill those gaps. A
\emph{per-frequency fill} admits any pixel exceeding a calibrated per-bin noise floor
$F[f]$ by \SI{2}{dB}, solidifying wideband bodies; a \emph{strong-signal rescue}
flags pixels \SI{8}{dB} above a re-estimated floor with a $\ge 3$-bin time-persistence
guard, recovering strong but frequency-narrow carriers.

\emph{Morphology and persistence.} The unioned mask is cleaned with a $3\times3$
majority filter, a $3\times7$ opening, and a $5\times21$ closing, followed by a
frequency-persistence test (keep a pixel only if $\ge 11$ of $45$ neighboring bins on
its time row are active); the strong-rescue mask is then OR-ed back to preserve
narrowband hits. A \SI{7}{MHz} sideband guard is applied per edge. The result is a
device-resident, full-frequency-resolution mask with population statistics attached as
metadata. The calibrated per-frequency floor makes this detector deterministic and
inexpensive, which is why it is the default operating point for real-time and
resource-constrained deployment; Table~\ref{tab:params} lists the active parameters.

\begin{table}[t]
\centering
\caption{Active coherent-power detector parameters (real-time single-channel
profile).}
\label{tab:params}
\scriptsize
\setlength{\tabcolsep}{2.5pt}
\begin{tabular}{@{}lll@{}}
\toprule
Stage & Parameter & Value \\
\midrule
Front end & noise smoothing $\sigma$ & 12 bins \\
          & max boost / signal cap & 0 dB (disabled) \\
Local background & time / frequency box radius & $\pm10$ / $\pm8$ bins \\
Support test & floor / span / threshold & 0.5 dB / 3 dB / 0.4 \\
Per-frequency fill & offset $\Delta_{\mathrm{fill}}$ & 2 dB \\
Strong rescue & excess / min.\ time bins & 8 dB / 3 \\
Morphology & opening / closing SE & $3\times7$ / $5\times21$ \\
           & majority filter & $3\times3$ \\
Persistence & window / minimum hits & 45 / 11 bins \\
Sideband guard & ignored per side & 7 MHz \\
\bottomrule
\end{tabular}
\end{table}

\subsection{Classical Baselines}
\label{sec:baselines}

Two non-learned detectors provide interpretable, low-cost reference points that run
directly on the dB-power frame. \emph{3dB\_power} is a single global threshold: a
per-frame scalar noise reference plus a fixed margin, flagging every bin above it.
\emph{blob\_detection} is a textbook computer-vision pipeline---Gaussian smoothing,
Sobel-gradient edges, percentile thresholding, morphological closing, hole filling,
and connected-component area filtering---deliberately left untuned. These establish
whether the coherent detector's added calibration, fill, and rescue machinery earns
its complexity over plain power thresholding and generic image processing.

\subsection{Learned Detectors}
\label{sec:learned}

Learned variants target the low-SNR, short-burst, and crowded regimes where fixed
power/morphology rules degrade. \emph{Zero-shot DINO} applies a frozen DINOv3
backbone to spectrogram tiles with no RF-specific training and runs as a live GPU
operator, testing whether generic self-supervised features transfer to RF as-is.
\emph{Fine-tuned DINOv3} adapts that backbone to RF: it trains a lightweight
segmentation head for dense signal/no-signal prediction \emph{and} fine-tunes the last
four transformer blocks and final norm of the backbone (at a reduced learning rate),
using captures spanning all attenuation levels. This head-plus-last-blocks adaptation
forms the high-recall operating point of the accuracy--compute knob. \emph{Fine-tuned YOLO} is a supervised object detector whose predicted boxes
are rasterized into the common mask, providing a strong discriminative-detector
comparison. All learned detectors emit the same mask interface and are scored
identically to the deployed detectors; Sec.~\ref{sec:eval} quantifies their
accuracy--compute tradeoff and their robustness under distribution shift.
